\documentclass[12pt]{article}
\usepackage[english]{babel}
\usepackage[utf8x]{inputenc}
\usepackage[T1]{fontenc}
\usepackage{mystyle}
\usepackage{csquotes}
\usepackage{listings}
\usepackage{hyperref}

% \DeclareNameAlias{sortname}{family-given}

% \defbibenvironment{bibliography}
%   {\begin{enumerate}}
%   {\end{enumerate}}
%   {\item}

\addbibresource{mybib.bib}



\Author{Sahana Rayan}
\Date{\today}
\Title{Literature Review: Neural Operators}

\lstset{style=mystyle}

\begin{document}
	\MakeScribeTop

\section{Background} \label{sec:background}

\subsection{Power Series} \label{sec:power}

Reference: \texttt{https://web.math.ucsb.edu/~moore/pde.pdf}.

Infinite Series of the following types is called a power series:
\begin{align*}
    e^x &= \sum_{n = 0}^{\infty} \frac{1}{n!}x^n \\
    \cos x &= \sum_{k = 0}^{\infty} (-1)^k \frac{1}{(2k)!}x^{2k} \\
    \sin x &= \sum_{k = 0}^{\infty} (-1)^k \frac{1}{(2k + 1)!}x^{2k + 1}
\end{align*}
 
Power series, more generally, is of the following form
\begin{align*}
    \sum_{n = 0}^{\infty} a_n(x - x_0)^{n}
\end{align*}
and they are best used for approximations.

For any power series, there exists $R$, which is a nonnegative real number or $\infty$, such that
\begin{enumerate}
    \item the power series converges when $|x - x_0| < R$
    \item and the power series diverges when $|x - x_0| > R$
\end{enumerate}
$R$ is called radius of convergence.

For example, $s_N = \sum_{n = 0}^{N} x^n = \frac{1 - x^N}{1 - x}$. $\lim_{N \to \infty} s_N = \frac{1}{1 - x}$ if $|x| < 1$ otherwise it diverges. Here, $R = 1$. More generally, $\sum_{n = 0}^{N} \left(\frac{x}{b}\right)^n$ has a radius of convergence of $R = b$

\textbf{Ratio test}:
The radius of convergence is given by:
\begin{align*}
    R = \lim_{n \to \infty} \frac{|a_n|}{|a_{n + 1}|}
\end{align*}
as long as the limit exists. 

So, for $\sum_{n = 0}^{N} \left(\frac{x}{b}\right)^n$, $a_n = \frac{1}{b^n}$, so $\frac{|a_n|}{|a_{n + 1}|} = b$ and so $R = b$.

For the power series of $e^x$, $a_n = \frac{1}{n!}$ and so, $\lim_{n \to \infty} \frac{|a_n|}{|a_{n + 1}|} = \lim_{n \to \infty}  \frac{(n + 1)!}{n!} = \lim_{n \to \infty}  n + 1 = \infty$. This means that the power series converges for all $x$. On the other hand, $\sum_{n = 0}^{\infty} n!x^n$, the ratio is $\frac{1}{n + 1}$ and that limit is $0$. Which means the power series doesn't converge for any non-zero $x$.

\textbf{Comparison Test:} Suppose that there are two power series:
\begin{align*}
    \sum_{n = 0}^{\infty} a_n(x - x_0)^{n} \\
    \sum_{n = 0}^{\infty} b_n(x - x_0)^{n}
\end{align*}
They have radii of convergence $R_1, R_2$ respectively. If $|a_n| \leq |b_n|$ for all $n$, then $R_1 \geq R_2$. If $|a_n| \geq |b_n|$ for all $n$, then $R_1 \leq R_2$. In other words, power series with smaller coefficients will result in larger radius of convergence.

For example, with the power series of $\cos x$, $|a_n| \geq \frac{1}{n!}$. We also know that $a_n = \frac{1}{n!}$ for the power series of $e^x$ whose radius of convergence is $\infty$. Therefore, $\cos x$ must have a radius of convergence at least as large as that of $e^x$. Therefore, it also converges for all values of $x$.

A function $f(x)$ is said to be \textbf{real analytic} at $x_0$ if there is a power series $\sum_{n = 0}^{\infty} a_n(x - x_0)^{n}$ about $x_0$ with positive radius of convergence $R$ such that
$f(x) = \sum_{n = 0}^{\infty} a_n(x - x_0)^{n}$, for $|x - x_0| < R$


For example, $e^x$ is real analytic at any $x_0$, $e^x = e^{x_0}e^{x - x_0} =e^{x_0} \sum_{n = 0}^{\infty} \frac{1}{n!}(x - x_0)^n = \sum_{n = 0}^{\infty} a_n(x - x_0)^{n}$ where $a_n = \frac{e^{x_0}}{n!}$. Similarly $f(x) = x^n$ is real analytic at $x_0$ because $x^n = (x - x_0 + x_0)^n = \sum_{i = 0}^n \frac{n!}{i!(n - i)!}x_0^{n - i}(x - x_0)^i$ in which $a_i = \frac{n!}{i!(n - i)!}x_0^{n - i}$ for $i \leq n$ and $a_i = 0$ for $i > n$. 

The sum and product of real analytic functions is real analytic. Any polynomial is analytic at any $x_0$. THe quotient of two polynomial with no common factors is analytic at $x_0$ if and only if $x_0$ is not a zero of the denominator polynomial. 

\subsection{Solving differential equations with power series}

The equation of simple harmonic motion is as follows:
\begin{align*}
    \frac{d^2 y}{d x^2} + y = 0
\end{align*}
we could assume that $y$ is a power series of some sorts or $y = \sum_{n= 0}^{\infty} a_n x^n$. if that is the case, 
$$\frac{dy}{dx} = \sum_{n = 1}^{\infty} na_nx^{n - 1}$$
and differentiating again yields
$$\frac{d^2 y}{d x^2} = \sum_{n = 2}^{\infty} n(n - 1)a_nx^{n - 2}$$
if $n = m + 2$ then
$$\frac{d^2 y}{d x^2} = \sum_{m = 0}^{\infty} (m + 2)(m + 1)a_{m + 2}x^{m}$$
$m$ is just a dummy variable so you can swap it out for $n$
$$\frac{d^2 y}{d x^2} = \sum_{n = 0}^{\infty} (n + 2)(n + 1)a_{n + 2}x^{n}$$




\end{document}